%    Latex Template for ACA 2015
%%%%%%%%%%%%%%%%%%%%%%%%%%%%%%%%%%%%%%%%
%    DO NOT CHANGE THE SETTINGS BELOW
%%%%%%%%%%%%%%%%%%%%%%%%%%%%%%%%%%%%%%%%
\documentclass[11pt]{article}
\usepackage{epsfig}
\usepackage{graphicx}
\usepackage[latin1]{inputenc}
\usepackage[T1]{fontenc}
\usepackage{mathptmx}
\usepackage{url}



\begin{document}



%%%%%%%%%%%%%%%%%%%%%%%%
%%%%%%%%%%%%%%%%%%%%%%%%%%%%%%%%%%%%%%%%
%    DO NOT CHANGE THE SETTINGS ABOVE
%%%%%%%%%%%%%%%%%%%%%%%%%%%%%%%%%%%%%%%%
%
\vspace{-0.8cm}
\begin{flushleft}
\Large \textbf{\noindent
%%% PLEASE INSERT THE TITLE OF YOUR TALK HERE %%%%%%%
Please Insert the Title of Your Talk Here}
%%%%%%%%%%%%%%%%%%%%%%%%%%%%%%%%%%%%%%%%%%%%%%%%%%%%%%%%
\\
%%%%%%%%%%%%%%%%%%%%%%%%%%%%%%
\vspace{0.5cm}
\normalsize
 %%%%%%             NAMES OF AUTHORS         %%%%%%
 %%%%%% NAME OF SPEAKER SHOULD BE UNDERLINED %%%%%%%%%
\normalsize{
 %%%%%% FOR EXAMPLE %%%%%
 A. Maths$^1$, B. Sciences$^2$, \underline{C. Lemma}$^3$, D. Technology$^4$
 %%%%%%%%%%%%%%%%%%%%%%%%
} \\
\vspace{5mm}
\textit{\footnotesize
 %%%%%% AFFILIATION OF AUTHORS %%%%%%
$^1$ University of Copenhagen, Denmark,
\{maths,social-sciences\}@ku.dk\\
$^2$ Wilfrid Laurier University, Waterloo, Canada,
natural-sciences@wlu.ca\\
$^3$ CSIRO, Canberra, Australia,
lemma@csiro.au\\
$^4$ Massachusetts Institute of Technology, Cambridge, USA,
technology@mit.edu\\
 %%%%%%%%%%%%%%%%%%%%%%%%%%%%%%%%%%%%%
}
\end{flushleft}

Use this format (without changing the settings) to prepare your  abstract. The abstracts are used for evaluation purposes. Therefore, it is important that you describe  your intended presentation concisely and clearly, including the objectives, status, methodology, results, and significance of your study.
 Abstracts may contain equations, figures and references (see Eqs.~(\ref{eq1}, \ref{eq2}) and Refs.~\cite{KR,WK,ABS}). {\bf The extended abstract length is limited to 5  pages}.  Equations are to be centered horizontally, with equation numbers aligned with the right margin.  Paragraph text is full-justified.  Do not explicitly begin by writing the word ``abstract''.  Equations should be properly punctuated.  For instance, as
\begin{equation}
y=3x^2\  \  \  \mbox{and}\  \  \  x>0, \label{eq1}
\end{equation}
we see that
\begin{equation}
z=\sqrt{y}=\sqrt{3}x. \label{eq2}
\end{equation}

%\begin{eqnarray}
%\lim_{x \to +\infty} {\rm Abstract}(x) = 1\ {\rm page}.
%\label{eq1}
%\end{eqnarray}
Your abstract should be prepared in \LaTeX\ according to the template provided  and then \textbf{converted to PDF for submission.}   Abstracts not complying with this template may not be included in the book of abstracts of the ACA 2015. Abstract submission should be done  via the session organizers, they will confirm your submission. If your abstract is accepted, you will be invited to present your talk.
The book of abstracts will be available to the participants prior to the conference and will provide the audience with more information
about the papers being presented at the ACA 2015 Workshop. At least one author must be registered for your talk to be included in the program of the conference.

%%%%%%%%%%%%%%%%%%%%%%%%%%%%%%%%%%%%%%%


%%%%%%%%%%%%%%%%%%%%%%%%%%%%%%%%%%%%%%%%
%% BIBLIOGRAPHY
%%%%%%%%%%%%%%%%%%%%%%%%%%%%%%%%%%%%%%%%
\begin{thebibliography}{00}
\addcontentsline{toc}{chapter}{References}
\setlength{\itemsep}{-1mm}
\footnotesize
\bibitem{KR} B. Keller and I. Reiten, \textit{Cluster-titled algebras are Gorenstein and stably Calabi-Yau}, Adv. Math. \textbf{211}, 1, pp. 123-151 (2007).
\bibitem{WK} S. Washito and K. Karabashi, \textit{Use This Template}, in \textit{Proceedings of the 5th
International Conference on Photon Interaction} (ICPI-2012), San Francisco, USA, ed. A. Chen, pp. 22-24 (2012).
\bibitem{ABS} A.B. Smith, \textit{Describing the style for the AMMCS-CAIMS-2015 book of abstracts}, in C.D. Science (ed.), General Book of Wisdom, 2\textsuperscript{nd} ed., University of Queensland, pp. 23-42 (2011).
\end{thebibliography}
\normalsize
\end{document} 
